\documentclass{article}

\usepackage{amsfonts}

\begin{document}

pag36

\hfill

\textbf{C02. Planar Linear Systems}

\hfill

In this chapter we begin the study of systems of differential equations. A system
of differential equations is a collection of n interrelated differential equations
of the form

\[
\begin{matrix}
 x'_1= f_1(t, x_1, x_2, ...,x_n) \\
x'_2= f_2(t, x_1, x_2, ...,x_n) \\
\vdots  \\
x'_n= f_n(t, x_1, x_2, ...,x_n)
\end{matrix}
\]

Here the functions $f_j$ are real-valued functions of the n+1 variables $x_1, x_2, ..., x_n$, and t . Unless otherwise specified, we will always assume that the $f_j$ are $C^{\infty}$
functions. This means that the partial derivatives of all orders of the fj exist
and are continuous.

To simplify notation, we will use vector notation:


\[
X = \begin{pmatrix}
x_1 \\
\vdots \\
x_n
\end{pmatrix} 
\]

We often write the vector X as $(x_1, ..., x_n)$ to save space.

Our system may then be written more concisely as

\[
X' = F(t, X)
\]

where

\[
F(t,X) = \begin{pmatrix}
f_1(t, x_1, ..., x_n) \\
\vdots \\
f_n(t, x_1, ..., x_n)
\end{pmatrix} 
\]

A solution of this system is then a function of the form $X(t) =
(x_1(t ), ..., x_n(t))$ that satisfies the equation, so that

\[
X'(t) = F(t, X(t))
\]

where $X'(t) = (x'_1(t), ..., x_n(t))$. Of course, at this stage, we have no guarantee
that there is such a solution, but we will begin to discuss this complicated
question in Section 2.7.

The system of equations is called autonomous if none of the $f_j$ depends on
t , so the system becomes X'= F(X). For most of the rest of this book we will
be concerned with autonomous systems.

In analogy with first-order differential equations, a vector X0 for which
$F(X_0) = 0$ is called an equilibrium point for the system. An equilibrium point
corresponds to a constant solution $X(t) \equiv X_0$ of the system as before.

Just to set some notation once and for all, we will always denote real variables
by lowercase letters such as $x, y, x_1, x_2, t$, and so forth. Real-valued functions
will also be written in lowercase such as f(x, y) or $f_1(x_1, ..., x_n, t)$. We will
reserve capital letters for vectors such as $X = (x_1, ..., x_n)$, or for vector-valued
functions such as

\[
F(x,y) = (f(x,y), g(x,y))
\]

or

\[
H(x_1, ..., x_n) = \begin{pmatrix}
h_1(x_1, ..., x_n) \\
\vdots \\
h_n(x_1, ..., x_n)
\end{pmatrix} 
\]

We will denote n-dimensional Euclidean space by $\mathbb{R}^n$, so that $\mathbb{R}^n$ consists of
all vectors of the form $X = (x_1, ..., x_n)$.

\hfill

\textbf{2.1 Second-Order Differential Equations}

\hfill

Many of the most important differential equations encountered in science
and engineering are second-order differential equations. These are differential
equations of the form

\[
x'' = f(t, x, x').
\]

Important examples of second-order equations include Newton’s equation

\[
mx''= f (x),
\]

the equation for an RLC circuit in electrical engineering

\[
LC x'' + RC x' + x = v(t),
\]

and the mainstay of most elementary differential equations courses, the forced
harmonic oscillator

\[
m x'' + b x' + kx = f (t).
\]

We will discuss these and more complicated relatives of these equations at length as we go along. First, however, we note that these equations are a special subclass of two-dimensional systems of differential equations that are defined by simply introducing a second variable y = x'.

For example, consider a second-order constant coefficient equation of the form

\[
x''+ax'+bx=0.
\]

If we let $y = x'$, then we may rewrite this equation as a system of first-order
equations

\[
\begin{matrix}
x' = y \\
y' = -bx - ay.
\end{matrix}
\]

Any second-order equation can be handled in a similar manner. Thus, for the
remainder of this book, we will deal primarily with systems of equations.

\hfill

\textbf{2.2 Planar Systems}

\hfill

For the remainder of this chapter we will deal with autonomous systems in
$\mathbb{R}^2$, which we will write in the form

\[
\begin{matrix}
x' = f(x,y) \\
y' = g(x,y).
\end{matrix}
\]

thus eliminating the annoying subscripts on the functions and variables. As
above, we often use the abbreviated notation X' = F(X) where X = (x, y)
and F(X) = F(x, y) = (f (x, y), g (x, y)).

In analogy with the slope fields of Chapter 1, we regard the right-hand side
of this equation as defining a vector field on $\mathbb{R}^2$. That is, we think of F(x, y)
as representing a vector whose x- and y-components are f (x, y) and g (x, y),
respectively. We visualize this vector as being based at the point (x, y). For
example, the vector field associated to the system

\[
\begin{matrix}
x' = y \\
y' = x.
\end{matrix}
\]

is displayed in Figure 2.1. Note that, in this case, many of the vectors overlap,
making the pattern difficult to visualize. For this reason, we always draw a
direction field instead, which consists of scaled versions of the vectors.
A solution of this system should now be thought of as a parameterized curve
in the plane of the form (x(t ), y(t )) such that, for each t , the tangent vector at
the point (x(t ), y(t )) is F(x(t ), y(t )). That is, the solution curve (x(t ), y(t ))
winds its way through the plane always tangent to the given vector F(x(t ), y(t ))
based at (x(t ), y(t )).

\hfill

\textbf{Example.} The curve

\[
\begin{pmatrix}
x(t) \\
y(t)
\end{pmatrix} = \begin{pmatrix}
a \sin t \\
a \cos t
\end{pmatrix}
\]

for any $a \in \mathbb{R}$ is a solution of the system

\[
\begin{pmatrix}
x' = y \\
y' = -x
\end{pmatrix}
\]

since

\end{document}